\paragraph{FastAPI: ``Default values'' (Jaccard 0.50).} \emph{0.70.0:} The same way that you can pass None as the first argument to be used as the default value, you can pass other values. Let's say that you want to declare the q query parameter to have a min\_length of 3, and to have a\,\dots{} \emph{0.115.0:} You can, of course, use default values other than None. Let's say that you want to declare the q query parameter to have a min\_length of 3, and to have a default value of "fixedquery": //// tab | Python 3.9+ //// //// tab | Python\,\dots

\paragraph{pandas: ``Inferring datetime format'' (Jaccard 0.50).} \emph{v1.5.3:} If you have parse\_dates enabled for some or all of your columns, and your datetime strings are all formatted the same way, you may get a large speed up by setting infer\_datetime\_format=True. If set, pandas will attempt to guess the format of your datetime strings,\,\dots{} \emph{v2.2.0:} Here are some examples of datetime strings that can be guessed (all representing December 30th, 2011 at 0000): * "20111230" * "2011/12/30" * "20111230 0000" * "12/30/2011 0000" * "30/Dec/2011 0000" * "30/December/2011 0000" Note that format inference is sensitive to dayfirst. With dayfirst=True, it\,\dots

\paragraph{scikit-learn: ``Locally Linear Embedding'' (Jaccard 0.51).} \emph{0.24.2:} Locally linear embedding (LLE) seeks a lower-dimensional projection of the data which preserves distances within local neighborhoods. It can be thought of as a series of local Principal Component Analyses which are globally compared to find the best non-linear embedding. Locally linear embedding can be\,\dots{} \emph{1.4.0:} Locally linear embedding (LLE) seeks a lower-dimensional projection of the data which preserves distances within local neighborhoods. It can be thought of as a series of local Principal Component Analyses which are globally compared to find the best non-linear embedding. Locally linear embedding can be\,\dots

\paragraph{Pydantic: ``Parsing data into a specified type'' (Jaccard 0.51).} \emph{v1.10.13:} Pydantic includes a standalone utility function parse\_obj\_as that can be used to apply the parsing logic used to populate pydantic models in a more ad-hoc way. This function behaves similarly to BaseModel.parse\_obj, but works with arbitrary pydantic-compatible types. This is especially useful when you want\,\dots{} \emph{v2.5.0:} [TypeAdapter][pydantic.type\_adapter.TypeAdapter] can be used to apply the parsing logic to populate Pydantic models in a more ad-hoc way. This function behaves similarly to [BaseModel.model\_validate][pydantic.main.BaseModel.model\_validate], but works with arbitrary Pydantic-compatible types. This is especially useful when you want to parse results into a type that is\,\dots
